% Generated draft: requires mathematical and source review.
\documentclass[11pt]{article}
\usepackage[margin=25mm]{geometry}
\usepackage{fontspec}
\setmainfont{Latin Modern Roman}
\usepackage{amsmath,amssymb,mathtools,unicode-math}
\setmathfont{Latin Modern Math}
\usepackage{enumitem,needspace,xurl,hyperref}
\setlength{\parindent}{0pt}
\setlength{\parskip}{5pt}
\setlength{\emergencystretch}{3em}
\newcommand{\N}{\mathbb{N}}
\newcommand{\Z}{\mathbb{Z}}
\newcommand{\Q}{\mathbb{Q}}
\newcommand{\R}{\mathbb{R}}
\newcommand{\C}{\mathbb{C}}
\newcommand{\F}{\mathbb{F}}
\newcommand{\A}{\mathbb{A}}
\newcommand{\PP}{\mathbb{P}}
\newcommand{\E}{\mathbb{E}}
\newcommand{\eps}{\varepsilon}
\newcommand{\dd}{\,\mathrm{d}}
\begin{document}
\section*{UnsolvedMath ID: 10000008}
\textbf{Source problem code:} catalogue entry
\subsection{Definitions and mathematical statement}
Let \(G=(V,E)\) be an infinite connected undirected graph with \(\sup_{v\in V}\deg_G(v)<\infty\), unit edge lengths and simple random walk. It is transient if the expected total number of visits to any fixed vertex is finite. It is Gromov-hyperbolic if some \(\delta<\infty\) makes every geodesic triangle \(\delta\)-thin in the graph metric.

The proposed assertion is that every such transient hyperbolic graph contains a transient subtree: there exists a connected acyclic subgraph \(T\subseteq G\) whose simple random walk is transient. The subtree need not be spanning or induced. Its transience concerns its own graph degrees and unit-edge random walk, not the walk on \(G\) restricted to its vertices.

The registry labels this assertion false in general and reports transient bounded-degree hyperbolic counterexamples with no transient subtree, even Liouville examples. Liouville means every bounded harmonic function on the graph is constant. This is a supplied resolution report without an identifying primary reference, not a proof or independently verified counterexample here. It conflicts with the provenance's bounded audit reporting no complete primary-literature resolution, so the original statement is retained together with the unresolved status evidence.
\subsection{Short English statement}
Must every bounded-degree transient Gromov-hyperbolic graph contain a subtree whose own simple random walk is transient? The registry reports counterexamples but supplies no reference and conflicts with its provenance audit.
\subsection{Sources}
\begin{itemize}
\item Local UnsolvedMath registry, permanent ID 10000008, source code catalogue entry; attributed to Itai Benjamini; no proposal year supplied; updated 2026-08-06; dated status solved.
\item Benjamini, \emph{Coarse Geometry and Randomness}, Saint-Flour notes, Open problem 1.62, PDF page 14: \url{https://www.wisdom.weizmann.ac.il/~itai/stflouraug24.pdf#page=14}. Supplied preserved copy: \url{https://arquivo.pt/noFrame/replay/20201231041548id_/http://www.wisdom.weizmann.ac.il/~itai/stflouraug24.pdf}. Registry access date 2026-07-29; text extraction checked against rendered PDF, as reported in provenance.
\item The source-side audit retains the item after deduplication, records no source-side resolution and no complete resolution in its bounded modern audit. The research summary nevertheless reports counterexamples, also relating them to a Benjamini--Schramm conjecture, without an author, title or citation. Both reports need reconciliation. Release reuse terms require review.
\end{itemize}
\end{document}
